\documentclass{article}
\usepackage[utf8]{inputenc}
\usepackage[spanish]{babel}
\usepackage[letterpaper, portrait, margin=2cm]{geometry}
\usepackage[style=ieee]{biblatex}
\usepackage{amsthm}
\usepackage{amsmath}
\usepackage{amssymb}
\usepackage{amsfonts}
\usepackage{hyperref}
\usepackage{csquotes}
\usepackage{mathtools}
\usepackage{graphicx}
\usepackage{float}
\usepackage{array}
\usepackage{dirtytalk}
\usepackage[table,xcdraw]{xcolor}
\usepackage{listings}

\addbibresource{bibliography.bib}

\newtheorem{definition}{Definición}

\DeclarePairedDelimiter{\ceil}{\lceil}{\rceil}

\definecolor{codegreen}{rgb}{0,0.6,0}
\definecolor{codegray}{rgb}{0.5,0.5,0.5}
\definecolor{codepurple}{rgb}{0.58,0,0.82}
\definecolor{backcolour}{rgb}{0.95,0.95,0.92}

\lstdefinestyle{mystyle}{
    backgroundcolor=\color{backcolour},   
    commentstyle=\color{codegreen},
    keywordstyle=\color{magenta},
    numberstyle=\tiny\color{codegray},
    stringstyle=\color{codepurple},
    basicstyle=\ttfamily\footnotesize,
    breakatwhitespace=false,         
    breaklines=true,                 
    captionpos=b,                    
    keepspaces=true,                 
    numbers=left,                    
    numbersep=5pt,                  
    showspaces=false,                
    showstringspaces=false,
    showtabs=false,                  
    tabsize=2
}

\lstset{style=mystyle}



\title{Proceso Propuesto de Investigación: Modelo de Muestreo Virtual de Suelos Agrícolas en Honduras para Analítica en Base al Muestreo Aleatorio Simple y Estratificado}
\author{Tobias Briones \bigbreak tobias.briones@unah.hn}
\date{Diciembre 2021}

\begin{document}

\makeatletter
    \begin{titlepage}
        \begin{center}
            \includegraphics[width=0.3\linewidth]{ref/logo-unah.png}\\[4ex]
            {\huge \bfseries \@title 
            \vspace{1cm}}\\[2ex]
            {\LARGE \@author}\\[50ex] 
            
            {\large
            Proyecto de Seminario de Investigación Presentado a la\\
            Universidad Nacional Autónoma de Honduras de la Carrera de\\
            Licenciatura en Matemática
            }\\[2ex]
            
            {\large \@date}
        \end{center}
    \end{titlepage}
\makeatother
\thispagestyle{empty}
\newpage

\section{Proceso Propuesto}

Se va a enunciar a continuación un problema particular del proceso propuesto determinado por este estudio. El problema es un producto mínimo viable que funciona como artefacto de investigación. El problema consiste en una exageración donde se toma el cultivo de azúcares en Honduras (literalmente todo el territorio de Honduras) y el terreno se clasifica por departamento. Para obtener los estratos, como se ha mencionado previamente, es conveniente recurrir a las variedades por lo que será útil definir algunas variedades de azúcares. Por último se ha de elaborar el muestreo virtual y reporte de los datos obtenidos.

\bigbreak

El analista de datos deberá interpretar y generalizar este proceso en una base de caso por caso para poder aplicarlo correctamente. La etapa de analítica es muy pesada de trabajo ya que hay que abrir y definir junto con los demás interesados todo el contexto del problema. En este ejemplo se trata con cultivos de azúcar y un terreno ya predispuesto así como la selección de azúcares por variedad como estrategia de homogeneización. Todas estas condiciones van a cambiar en un caso real pero se pueden considerar abstraídas mediante estos resultados por lo que el analista de datos podrá ser capaz de generalizar.

\bigbreak

Los datos en particular se generarán de forma aleatoria al no contar con un conjunto de datos reales ya que esto estaría fuera del alcance de investigación.

\bigbreak

En los párrafos con un símbolo de implicación se dará una generalización del párrafo anterior para el analista de datos.

\subsection{Enunciado}

Asumir que en todo el territorio de Honduras se cultiva azúcar, es decir --la finca o suelo agrícola consiste en todo el territorio del país-- y se proporcionan los datos de Zafra 2021. Mediante un SIG se ha obtenido un mapa del área clasificada por departamento, en total, 18 departamentos.

\bigbreak

$\implies$ El analista de datos en esta etapa deberá contar con la muestra física y el modelo del suelo agrícola probablemente en formato shapefile (shp). Por tanto, el conjunto de datos puede comprender a cualquier tipo de  cultivo además de azúcar y cualquier zona geográfica que deberá tener
cierta delimitación geográfica.

\bigbreak

A fin de obtener los estratos se ha establecido usar la variedad --variable que define propiedades biológicas de los cultivos-- para clasificar el suelo. Al ser cosecha de azúcar, se asume que se cultivan las variedades:

\begin{itemize}
    \item Ninguna
    \item CP 72-2086
    \item Mex 69-290
    \item Mex 79-431
    \item ITV 92-1424
\end{itemize}

Las cuales son variedades muy conocidas.

\bigbreak

$\implies$ El analista obtendrá una lista mucho más grande de variedades disponibles en el suelo.

\bigbreak

Se debe de obtener un muestreo físico con la siguiente estructura mínima:

\begin{itemize}
    \item Estrato
    \item Área Cosechada (HA)
    \item TA/HA Real
\end{itemize}

Donde el conjunto de datos representa los registros tomados por el usuario en ese año de zafra.

\subsection{Análisis Geográfico}

Este proceso consiste en un análisis del SIG estándar. Se definen las constantes, se carga el archivo shapefile y se le da formato correspondiente agregando demás columnas. Se identifica como variable independiente a la columna departamento.

\bigbreak

El análisis del SIG da un reporte de alto nivel sobre el muestreo físico. Luego se carga el \textit{DataFrame} del muestreo físico para zafra año 2021 en este caso, entonces se desarrolla el MAS Estratificado y se establece el reporte del muestreo obtenido. Con esto, se puede hacer finalmente un análisis de fertilidad consistente y más eficiente. Por tanto, el análisis geográfico es la de las tres etapas mencionadas.

\subsubsection{Recopilación de Datos}

Este problema es una exageración que permite de momento llevar a cabo el proceso propuesto de muestreo virtual. La mayor parte de datos creados en este problema no tienen relación con la realidad pero si sirven como modelo para ser aplicado a un caso con datos reales.

\bigbreak

Los datos del muestreo físico son recolectados mediante la desnormalización de la base de datos del usuario. Esto implica que en análisis de datos se ven muchos fenómenos como redundancia, etc. El cultivo de azúcares se da por zafra, en este caso, se asume una zafra 2021 la cual da todos los registros de ese año.

\bigbreak

Las 4 \href{https://www.gob.mx/conadesuca/es/articulos/variedades-de-cana-de-azucar?idiom=es}{variedades} tomadas son una referencia importante que se encontrará de forma mucho más diversa en una base de datos real. Estas son variedades significativas, en un caso real se tiene una gran cantidad de variedades de azúcares por lo que el analista deberá tomar solo aquellas significativas para reducir la carga de la muestra. Otras variedades de azúcares significativas incluye, a saber, \textit{CP 73-1547}, \textit{RB 86-7515}, etc. Las encuestas pueden ser herramientas adecuadas para la obtención de esta información por parte de los interesados.

\bigbreak

Por tanto, se debe de tomar una variable adecuada física o biológica como la variedad del cultivo para crear cada estrato con unidades de muestreo homogéneas. Al definir inicialmente los estratos, se tendrá una lista muy grande de opciones por lo que será necesario tomar un subconjunto significativo de estos estratos filtrando aquellos que no son significativos de acuerdo al criterio establecido por el analista. Después de este filtrado, se obtienen los estratos definitivos y el mapa original queda con áreas no representativas por lo que se define el estrato vacio \textit{Ninguno}. Originalmente puede que ya existan áreas no representativas como zonas de construcción, contaminadas, etc.

\subsubsection{Filtrado}

El filtrado inicial se ha definido como sigue:

\begin{itemize}
    \item Remover estratos vacíos.
    \item Remover estratos con variedades insignificantes.
\end{itemize}

Para determinar las variedades insignificantes se establece una cota mínima (umbral) de área cosechada por variedad, estos datos vienen del muestreo físico. Por tanto, la condición para determinar si una variedad es significante es si pertenece a las cuatro variedades definidas arriba o si el área cosechada para esa variedad supera un umbral determinado. En este problema solo se filtrará tomando las cuatro variedades comúnes.

\subsection{Muestreo Físico}

Este \textit{DataFrame} contiene los supuestos datos de zafra $2021$ en el suelo agrícola. Es decir, todos los registros de variedades tomados en ese año. Además, en un conjunto de datos original existen muchas otras variables que pueden ser usadas para otros cálculos. Este \textit{DataFrame} solo contiene el filtrado con las columnas correspondientes para el análisis de rendimiento.

\subsection{Resumen de Muestro}

En las siguientes dos tablas se listan los datos del SIG y del muestreo físico filtrado para análisis de rendimiento correspondientemente.

\subsubsection{Modelo de Muestreo}

\textbf{Muestro Físico}

\bigbreak

El tamaño de la población $N = 5,000$ unidades de muestreo.

\bigbreak

Unidad de área: Hectárea.

\bigbreak

\textbf{Muestreo Virtual}

\bigbreak

\textbf{Estratos}

\bigbreak

Se han definido cuatro estratos significativos correspondientes a variedades de azúcares y uno vacío. Por tanto, $H = \{ Ninguno, H_1, ... , H_4 \}$.

\bigbreak

Tamaño de muestreo por estrato $(H_n) = t_h * factor$ donde \textit{factor} en $[0,1]$ representa el porcentaje de unidades aleatorias a tomar en cada estrato $H$ respecto a $t_h$.

\bigbreak

Con respecto al conjuntos de datos geográficos originalmente cargado, este simplemente cuenta con una (generosa) columna \textit{geography} la cual contiene la información vectorial del mapa del suelo agrícola. También fue útil la columna \textit{DEPARTMENT\_COL} conteniendo la partición hecha por el usuario. En este caso, los estratos son por variedad significante por lo que el problema debe ser tal que los subconjuntos de la partición del mapa (\say{departamentos}) contienen únicamente una variedad. Esto se detalla más en la sección de visualización abajo. A partir de estos criterios, el analista deberá poder obtener la información geográfica útil del usuario.

\subsection{Visualización de Muestro}

En los siguientes bloques se visualizan los datos obtenidos arriba.

\bigbreak

Por ejemplo, en los \say{departamentos} en azul (Cortés, Atlántida, Yoro,Ocotepeque, Lempira y El Paraíso) se cultiva variedad \textit{CP 72-2086}. Así, esos $6$ \say{departamentos} forman un estrato.


\printbibliography

\end{document}
